\chapter{Introduction}
\label{ch:introduction}

\section{Contexte}
\label{sec:contexte}

Les bras robotiques low-cost (ici, le SO-101 de Hugging Face, dérivé du
SO-ARM100 de TheRobotStudio) sont de plus en plus accessibles, mais restent
limités dès qu'ils doivent interagir de manière autonome avec des
environnements réels et encombrés. En particulier, lorsqu'un objet cible est
entouré d'autres objets, le robot ne peut pas se contenter d'une trajectoire
prédéfinie ni d'une stratégie de saisie unique~\cite{bohg_data-driven_2014}.

La perception visuelle est alors sujette à des problèmes d'occlusion, de
points de vue incomplets et d'incertitudes sur la géométrie de l'objet. De
plus, le chemin le plus court vers l'objet cible n'est pas toujours exploitable
en raison des risques de collision avec les distracteurs présents dans la
scène.

\section{Problématique}
\label{sec:problematique}

La problématique centrale de ce travail s'énonce ainsi : \emph{comment décider
dynamiquement où regarder et comment saisir un objet dans un environnement
encombré, en combinant plusieurs sources de perception visuelle, tout en
garantissant l'évitement d'obstacles et une exécution robuste du mouvement ?}

Ce travail s'inscrit dans une approche \textbf{perception--action}~\cite{ghallab_lanticipation_2015},
où la difficulté ne réside pas uniquement dans la calibration ou la détection
de l'objet, mais dans la prise de décision et l'orchestration entre perception,
planification de trajectoire et contrôle du bras robotique.

\section{Objectifs}
\label{sec:objectifs}

L'objectif principal de ce travail de fin d'études est de concevoir,
implémenter et évaluer une architecture permettant à un bras robotique équipé
de caméras eye-to-hand et eye-in-hand d'interagir de manière autonome avec un
environnement encombré.

Plus précisément, les objectifs sont :
\begin{enumerate}
  \item Concevoir une architecture
        \textbf{perception--planification--contrôle} combinant les informations
        issues de plusieurs caméras.
  \item Proposer une stratégie de \textbf{sélection de points de vue} afin
        d'améliorer la perception de l'objet cible en présence d'occlusions.
  \item \textbf{Planifier} des trajectoires de mouvement évitant les obstacles
        présents dans l'environnement.
  \item Mettre en œuvre une \textbf{replanification en boucle fermée} basée
        sur la perception, permettant d'adapter la trajectoire du robot en
        fonction des informations visuelles mises à jour au cours du mouvement.
  \item Définir et sélectionner des stratégies de \textbf{saisie} adaptées à
        la géométrie et à l'accessibilité de l'objet (grasp planning).
  \item \textbf{Évaluer expérimentalement} le système proposé à l'aide de
        métriques telles que le taux de réussite de la saisie, le nombre de
        replans, les collisions évitées et la précision de la prise.
\end{enumerate}

Dans un premier temps, la prise de décision repose sur des règles heuristiques
et des critères géométriques (chapitre~\ref{ch:perception} pour la perception,
chapitre~\ref{ch:planning} pour la planification de saisie). Des approches
plus avancées, telles que l'apprentissage à partir de démonstrations
(\textit{imitation learning}, abordé au chapitre~\ref{ch:control_eval}), sont
envisagées comme extension comparative dans la stack officielle
LeRobot~\cite{shukor_smolvla_2025}.

\section{Matériel}
\label{sec:materiel}

Le poste de travail expérimental est constitué (voir
figure~\ref{fig:setup}~--- \todo{ajouter photo}) :

\begin{itemize}
  \item \textbf{Robot} : SO-101 (Hugging Face / TheRobotStudio), 6 servomoteurs
        Feetech STS3215 par bras (leader + follower), gripper à pince parallèle.
  \item \textbf{Caméras} : 3 caméras USB 1920$\times$1080 à 30~FPS :
    \begin{itemize}
      \item \texttt{cam\_0} et \texttt{cam\_1} : configuration
            \textbf{eye-to-hand stéréo} sur la barrière avant du poste.
      \item \texttt{cam\_2} : configuration \textbf{eye-in-hand} montée sur
            l'effecteur du bras.
    \end{itemize}
  \item \textbf{Structure 3D} : poste de travail modélisé sous CAO, imprimé en
        3D, et assemblé définitivement. Inclut un point de dépose fixe (boîte
        de $15 \times 10 \times 6$ cm).
  \item \textbf{Machine de développement} : MacBook Pro M4, 24~GB RAM, macOS.
\end{itemize}

\begin{figure}[hbtp]
\centering
\rule{0.6\textwidth}{0.5\textwidth}
\caption{Vue d'ensemble du poste expérimental SO-101 avec les trois caméras.
\todoF{Remplacer par une vraie photo du poste.}}
\label{fig:setup}
\end{figure}

\section{Plan du mémoire}
\label{sec:plan}

Le chapitre~\ref{ch:etat_de_lart} dresse l'état de l'art sur la perception
multi-caméras, la triangulation stéréo, les détecteurs d'objets modernes
(seuillage classique vs détecteurs open-vocabulary), et les deux grands
paradigmes contemporains de la manipulation robotique
(\textit{architecture modulaire} vs \textit{imitation learning end-to-end}).
Le chapitre~\ref{ch:architecture} présente l'architecture logicielle adoptée.
Les chapitres~\ref{ch:perception} à \ref{ch:control_eval} détaillent
l'implémentation des modules de perception, de planification de saisie et de
contrôle. Le chapitre~\ref{ch:discussion} discute les limites identifiées et
les pistes d'amélioration. Le chapitre~\ref{ch:conclusion} conclut.
